\documentclass{article}
\usepackage[T1]{fontenc}
\usepackage{lmodern}
\usepackage{graphicx}
\usepackage{amsmath,amssymb}
\usepackage[a4paper,margin=2cm]{geometry}
\usepackage[absolute,overlay]{textpos}
\setlength{\TPHorizModule}{1mm}
\setlength{\TPVertModule}{1mm}
\usepackage[indonesian]{babel}
\usepackage{hyperref}
\setlength{\emergencystretch}{2em}
% unit: Halaman 21

\begin{document}

% === Halaman 21 (python/native) ===

Implementasi Knapsack

• Ukuran cipherteks yang dihasilkan lebih besar daripada plainteksnya, 

karena enkripsi dapat menghasilkan kriptogram yang nilai desimalnya lebih 

besar daripada nilai desimal blok plainteks yang dienkripsikan.

• Untuk menambah kekuatan algoritma knapsack, kunci publik maupun 

kunci privat seharusnya paling sedikit 250 elemen, nilai setiap elemen 

antara 200 sampai 400 bit panjangnya, nilai modulus antara 100 sampai 

200 bit. 

• Dengan nilai-nilai knapsack sepanjang itu, dibutuhkan 1046 tahun untuk 

menemukan kunci secara brute force, dengan asumsi satu juta percobaan 

setiap detik.

21

\end{document}
